\chapter*{Wprowadzenie do tomu II}
\addcontentsline{toc}{chapter}{Wprowadzenie do tomu II}
Tom I dotyczył zmiennych losowych, ich rozkładów, zbieżności, warunkowej
wartości oczekiwanej, łańcuchów Markowa i martyngałów w czasie dyskretnym.
Tom II bada całe \emph{rodziny} zmiennych losowych, indeksowane czasem,
w szczególności czasem ciągłym. Jest to zmiana perspektywy: oprócz rozkładu
w danej chwili interesuje nas wspólna zależność wielu chwil, a także
regularność całej trajektorii.

W pierwszej części odpowiemy na trzy odrębne pytania. (1) Czym dokładnie
jest proces i co znaczy znać jego rozkład? (2) Czy zadanie zgodnych
rozkładów skończeniewymiarowych wyznacza rzeczywisty proces, a jeśli tak,
kiedy ma on ciągłe trajektorie? (3) Jak formalnie opisać narastającą
w czasie informację i podejmować decyzje bez znajomości przyszłości?
Rozdziały zostały ułożone tak, by konstrukcja procesu Wienera,
miary Wienera i regularności trajektorii opierała się na ścisłych
fundamentach probabilistycznych. Całka stochastyczna Itô,
równania stochastyczne i systematyczna teoria całek po trajektoriach
będą rozwinięte w tomie III, poświęconym analizie stochastycznej.
Tom IV obejmie statystykę matematyczną i wnioskowanie bayesowskie.

Zakładamy znajomość teorii miary i całki Lebesgue'a, twierdzenia
Radona--Nikodyma, warunkowej wartości oczekiwanej, lematu
Borela--Cantellego oraz podstaw zbieżności zmiennych losowych z tomu I.
Jako znane wyniki z teorii miary traktujemy twierdzenie Carathéodory'ego
o rozszerzaniu miar, regularność borelowskich miar probabilistycznych
na przestrzeniach polskich oraz twierdzenie o zwartości iloczynu
przeliczalnej rodziny niepustych przestrzeni zwartych.
\emph{Nie} zakładamy twierdzenia Kołmogorowa o istnieniu procesu ani jego
kryterium ciągłości: oba zostaną udowodnione.

Równości i nierówności między zmiennymi losowymi rozumiemy
jako zachodzące prawie na pewno, o ile nie zaznaczono inaczej.
W rozdziałach 1--3 podkreślimy, że pojęcie modyfikacji nie wymaga
równości procesów poza jednym zbiorem zerowym dla wszystkich chwil.
To odróżnia je od nierozróżnialności.

